\begin{center}
\LARGE
Groebner Basis: A Bridge between Coding Theory and Computer Algebra
\end{center}

\begin{center}
\Large
Shojiro Sakata
\end{center}

\noindent
Date:  July 17th (Wednesday)\\
Time:  16:20-17:00 \\

\begin{center}
\large
Abstract
\end{center}

In this talk we give a survey of 
\begin{enumerate}
\item
when the concept of Groebner basis has been introduced into the world of coding theory; 
\item
how the concept of Groebner basis has been applied to construction and decoding of
    error-correcting codes,
\end{enumerate}
and discuss some important roles of Groebner basis in recent development of coding theory. It is
emphasized that Buchberger's algorithm and its relatives take their parts just at the point of contact
of coding theory with system theory and that Groebner basis is a key not only to full decoding of
the most frequently used error-correcting codes such as Reed-Solomon and
Bose-Chaudhuri-Hocquenghem codes but also to construction and efficient decoding method of
the next generation of error-correcting codes called algebraic-geometry codes. 


